\documentclass[12pt]{article}
\usepackage[T2A]{fontenc}
\usepackage[utf8]{inputenc}
\usepackage[ukrainian]{babel}
\usepackage[hidelinks]{hyperref}
\usepackage{amsmath,amssymb,amsthm,thmtools}
\usepackage{mathtools}
\usepackage{enumitem,booktabs}
\mathtoolsset{centercolon,mathic}
\usepackage{listings}
\usepackage{xcolor,xspace}
\usepackage{tikz-cd}
\usetikzlibrary{babel,positioning,arrows.meta,decorations.pathreplacing}
\usepackage{microtype}
\usepackage{url}

%%% TikZ setup
\tikzset{thick red/.style={thick,color=darkred}}
\tikzset{edgelabel/.style = {midway,font=\small}}
\tikzset{slopedlabel/.style = {edgelabel,sloped}}
\definecolor{darkred}{RGB}{170,0,0}
\tikzcdset{arrow style=tikz}

\tikzset{
  symbol/.style={
    draw=none,
    every to/.append style={
      edge node={node [sloped, allow upside down, auto=false]{$#1$}}
    }
  }
}

%%% Lean listing setup
\def\lstlanguagefiles{lstlean.tex}
\lstset{language=lean}
\definecolor{DarkBlue}{rgb}{0.03, 0.21, 0.48}

%%% MACROS
\newcommand*{\Loop}{\Omega}
\newcommand*{\concat}{\cdot}
\mathchardef\texthyphen="2D
\newcommand*{\susp}[1]{{\Sigma}{#1}}
\newcommand*{\sg}[1]{(#1) \times }
\newcommand*{\rev}[1]{{#1}^{-1}}
\newcommand*{\bN}{{\mathbb{N}}}
\newcommand*{\bF}{{\mathbb{F}}}
\newcommand*{\bZ}{{\mathbb{Z}}}
\newcommand*{\bR}{{\mathbb{R}}}
\newcommand*{\bC}{{\mathbb{C}}}
\newcommand*{\bH}{{\mathbb{H}}}
\newcommand*{\bS}{{\mathbb{S}}}
\newcommand*{\RP}{{\mathbb{R}\mathrm{P}}}
\newcommand*{\CP}{{\mathbb{C}\mathrm{P}}}
\newcommand*{\HP}{{\mathbb{H}\mathrm{P}}}
\newcommand*{\pt}{{\mathrm{pt}}}
\newcommand*{\limit}{\varprojlim}
\newcommand*{\colimit}{\varinjlim}
\newcommand*{\dblslash}{\mathbin{/\kern-3pt/}}
\renewcommand*{\equiv}{\mathrel{\simeq}}
\DeclarePairedDelimiter\abs{\lvert}{\rvert}
\DeclarePairedDelimiter\angled{\langle}{\rangle}
\DeclarePairedDelimiter\trunc{\lVert}{\rVert}
\DeclarePairedDelimiter\tr{\lvert}{\rvert}
\newcommand*{\pair}[2]{\angled{#1,#2}}
\DeclareMathOperator\Th{Th}
\DeclareMathOperator\Decat{Decat}
\DeclareMathOperator\Disc{Disc}
\DeclareMathOperator\aut{Aut}
\DeclareMathOperator\Aut{Aut}
\DeclareMathOperator\BAut{BAut}
\DeclareMathOperator\ap{ap}
\DeclareMathOperator\B{B}
\newcommand*\BS{\mathop{BS}}
\newcommand*\Sym{S}
\newcommand*\BZ{\mathop{\mathrm{B}\bZ}}
\DeclareMathOperator\id{id}
\DeclareMathOperator\im{im}
\DeclareMathOperator\inl{inl}
\DeclareMathOperator\inr{inr}
\DeclareMathOperator\fib{fib}
\DeclareMathOperator\Fin{Fin}
\DeclareMathOperator\istrunc{istrunc}
\DeclareMathOperator\iscontr{iscontr}
\DeclareMathOperator\isprop{isprop}
\DeclareMathOperator\isconn{isconn}
\newcommand*{\Prop}{\mathrm{Prop}}
\newcommand*{\Set}{\mathrm{Set}}
\newcommand*{\Type}{\mathrm{Type}}
\newcommand*{\GType}{\mathrm{GType}}
\newcommand*{\Grp}{\mathrm{\texthyphen Group}}
\newcommand*{\Group}{\mathrm{Group}}
\newcommand*{\AbGroup}{\mathrm{AbGroup}}
\newcommand*{\iDecat}{\mathrm{\infty\texthyphen\mathrm{Decat}}}
\newcommand*{\iDisc}{\mathrm{\infty\texthyphen\mathrm{Disc}}}
\renewcommand*{\phi}{\varphi}
\newcommand*{\blank}{{-}}
\newcommand*{\leanref}[1]{{\normalfont[\texttt{\color{DarkBlue}#1}]}}
\newcommand*{\ditto}{--- \raisebox{-0.5ex}{''} ---}

%%% THEOREMS
\declaretheorem[style=plain,numberwithin=section,name=Теорема]{theorem}
\declaretheorem[style=plain,sibling=theorem,name=Лема]{lemma}
\declaretheorem[style=plain,sibling=theorem,name=Пропозиція]{proposition}
\declaretheorem[style=plain,sibling=theorem,name=Наслідок]{corollary}
\declaretheorem[style=plain,sibling=theorem,name=Гіпотеза]{conjecture}
\declaretheorem[style=definition,sibling=theorem,name=Визначення]{definition}

\title{Вищі групи в гомотопічній теорії типів}
\author{Ульрік Бухгольц, Флоріс ван Доорн та Егберт Райке}
\date{\today}

\begin{document}

\maketitle

\begin{abstract}
  Ми представляємо розвиток теорії вищих груп, включаючи нескінченні групи та зв'язні спектри, у гомотопічній теорії типів. Нескінченна група — це просто тип петель у точковому зв'язному типі, де групова структура випливає зі структури, притаманної типам тотожності теорії типів Мартіна-Льофа. Ми досліджуємо звичайні групи з цієї точки зору, а також вищі розмірні групи та групи, які можна розпетлювати більше одного разу. Основним результатом є теорема про стабілізацію, яка стверджує, що якщо $n$-тип можна розпетлювати $n+2$ рази, то він є типом нескінченних петель. Більшість результатів було формалізовано у доказовому асистенті Lean.
\end{abstract}

\section{Вступ}
\label{sec:introduction}

\begin{table*}
  \caption{\label{tab:periodic}Періодична таблиця $k$-кратно групових $n$-групоїдів.}
\makebox[\textwidth][c]{
  \begin{tabular}{clllll} \toprule
    $k\setminus n$ & $0$ & $1$ & $2$ & $\cdots$ & $\infty$ \\
    \midrule
    $0$ & точковий набір & точковий групоїд & точковий $2$-групоїд & $\cdots$ & точковий $\infty$-групоїд \\
    $1$ & група & $2$-група & $3$-група & $\cdots$ & $\infty$-група \\
    $2$ & абелева група & плетена $2$-група & плетена $3$-група & $\cdots$ & плетена $\infty$-група \\
    $3$ & \ditto & симетрична $2$-група & силептична $3$-група & $\cdots$ & силептична $\infty$-група \\
    $4$ & \ditto & \ditto & симетрична $3$-група & $\cdots$ & ?? $\infty$-група \\
    $\vdots$ & \mbox{}\quad$\vdots$ & \mbox{}\quad$\vdots$ & \mbox{}\quad$\vdots$ & $\ddots$ & \mbox{}\quad$\vdots$ \\
    $\omega$ & \ditto & \ditto & \ditto & $\cdots$ & зв'язний спектр \\
    \bottomrule
  \end{tabular}%
  }
\end{table*}

Гомотопічна гіпотеза — це твердження про те, що гомотопічні $n$-типи (топологічні простори з тривіальними гомотопічними групами вище рівня $n$) відповідають $n$-групоїдам для $n \in \bN \cup \{\infty\}$ через конструкцію фундаментального $\infty$-групоїда. У оригінальній версії Гротендіка в \emph{Pursuing Stacks}~\cite{Grothendieck1983} це була гіпотеза про певну модель $\infty$-групоїдів. Це також є теоремою для багатьох конкрених моделей $\infty$-групоїдів, наприклад, симпліціальних множин Кана, але зараз це здебільшого сприймається як властивість, що \emph{визначає} $\infty$-групоїди з точністю до еквівалентності.

У цій статті ми досліджуємо гомотопічну гіпотезу в контексті гомотопічної теорії типів (HoTT). HoTT відноситься до гомотопічної інтерпретації залежної теорії типів Мартіна-Льофа~\cite{AwodeyWarren2009,Voevodsky06}. У цій гомотопічній інтерпретації кожна теоретико-типова конструкція відповідає гомотопічно інваріантній конструкції на просторах.

У HoTT кожен тип має простір шляхів, заданий типом тотожності. Для точкового типу ми можемо побудувати простір петель, який має структуру $\infty$-групи. Більше того, якщо тип \emph{усічений}, то ми можемо повернути звичайне поняття груп, $2$-груп та вищих груп. Це дозволяє нам визначити вищу групу внутрішньо мовою теорії типів як тип, що є простором петель точкового зв'язного типу, його розпетлювання.

Ми також досліджуємо групи, які можна розпетлювати більше одного разу, що дає $n$-групи з додатковими когерентностями. Повне сімейство груп, які ми розглядаємо, наведено в~\autoref{tab:periodic}, що ми детально пояснимо в~\autoref{sec:higher-groups}.

Наш підхід додатково підтверджується відповідним спостереженням у теорії $\infty$-топосів, де є теоремою те, що $\infty$-категорія точкових зв'язних об'єктів в $\mathcal X$ еквівалентна $\infty$-категорії вищих групових об'єктів в $\mathcal X$ для будь-якого $\infty$-топосу $\mathcal X$~\cite[Лема~7.2.2.11(1)]{LurieHTT}.

Ми формалізували більшість наших результатів у бібліотеці HoTT~\cite{vDvRB2017HoTTLean} доказового асистента Lean~\cite{Moura2015}. Формалізовані результати можна знайти у файлі \url{https://github.com/cmu-phil/Spectral/blob/master/higher_groups.hlean}. Ми будемо вказувати основні формалізовані результати в цій статті, посилаючись на назву у формалізації у квадратних дужках. Для отримання додаткової інформації про формалізацію див.~\autoref{sec:formalization}.

Ми вдячні Майклу Шульману за написання допису в блозі~\cite{Shulman-classifying} про класифікуючі простори з унівалентної перспективи.

\section{Попередні відомості}
\label{sec:preliminaries}

У цій статті ми будемо працювати в теорії типів книги HoTT~\cite{TheBook}, хоча всі аргументи також будуть справедливі в кубічній теорії типів, такій як~\cite{CCHM2016,chtt}. У цьому розділі ми коротко вводимо концепції, необхідні для решти статті.

Теорія типів містить залежні типи функцій $(x : A) \to B(x)$, які традиційно позначаються як $\Pi_{x : A}B(x)$, і залежні типи пар $(x : A) \times B(x)$, які традиційно позначаються як $\Sigma_{x : A}B(x)$. Ми вибираємо цю нотацію, натхненну Agda, тому що часто маємо справу з глибоко вкладеними типами залежних сум.

Всередині типу $A$ ми маємо тип тотожності або тип шляху ${=}_A : A \to A \to \Type$. Ми маємо різні операції над шляхами, такі як конкатенація $p \concat q$ та інверсія $\rev{p}$ шляхів. Функторіальна дія функції $f : A \to B$ на шлях $p : a_1 =_A a_2$ позначається $\ap_f(p) : f(a_1) = f(a_2)$. Постійний шлях позначається $1_a : a = a$.

Тип $A$ може бути $n$-усіченим, що позначається $\istrunc_n A$ і визначається рекурсією по $n : \bN_{-2} := \mathbb{Z}_{\ge-2}$:
\begin{align*}
  \istrunc_{-2} A & := \iscontr A := (a : A) \times \bigl((x : A) \to (a = x)\bigr) \\
  \istrunc_{n+1} A & := (x\;y : A) \to \istrunc_n (x=y)
\end{align*}
Для будь-вого типу $A$ ми пишемо $\trunc{A}_n$ для його $n$-усічення, тобто $\trunc{A}_n$ — це $n$-усічений тип, оснащений відображенням $\tr{\blank}_n : A \to \trunc{A}_n$, таким що для будь-якого $n$-усіченого типу $B$ відображення прекомпозиції
\begin{equation*}
(\trunc{A}_n \to B) \to (A \to B)
\end{equation*}
є еквівалентністю. Тоді ми визначаємо $n$-зв'язність як $\isconn_n A := \iscontr \trunc A_n$. Властивості усічень та зв'язних відображень встановлені у Розділі 7~\cite{TheBook}.

Тип точкових типів — це $\Type_\pt := (A : \Type) \times (\pt : A)$. Тип $n$-усічених типів — це $\Type^{\le n} := (A : \Type) \times \istrunc_n A$, а для $n$-зв'язних типів — це $\Type^{>n} := (A : \Type) \times \isconn_n A$. Ми будемо комбінувати ці позначення за потреби.

Для $A : \Type_\pt$ ми визначаємо \emph{простір петель} $\Loop A := (\pt =_A \pt)$, який є точковим з базовою точкою $1_\pt$. \emph{Гомотопічні групи} $A$ визначаються як $\pi_k A := \trunc{\Loop^k A}_0$. Вони є групами у звичайному розумінні при $k \ge 1$, з нейтральним елементом $\tr{1}$ і груповою операцією, індукованою конкатенацією шляхів.

Для $A, B : \Type_\pt$ тип \emph{точкових відображень} від $A$ до $B$ — це $(A \to_\pt B) := (f : A \to B) \times (f(\pt) =_B \pt)$. Для $f : A \to_\pt B$ ми пишемо $f(a) : B$ для першої проекції та $f_0 : f(\pt) = \pt$ для другої проекції. \emph{Волокно} точкового відображення визначається як $\fib(f) := (a : A) \times (f(a) =_B \pt)$, яке є точковим з базовою точкою $(\pt, f_0)$.

У HoTT ми можемо використовувати \emph{вищі індуктивні типи} для побудови просторів Ейленберга-Маклейна $K(G,n)$~\cite{FinsterLicata2014}. Для групи $G$ ми визначаємо $K(G,1)$ як наступний ВІТ.\\
\indent \texttt{HIT} $K(G,1) :=$ \\
\indent $\bullet\ \star : K(G,1)$; \\
\indent $\bullet\ p : G \to \star = \star$; \\
\indent $\bullet\ q : (g\;h : G) \to p(gh) = p(g) \concat p(h)$; \\
\indent $\bullet\ \epsilon : \istrunc_1 K(G,1)$. \\
(Використовуючи унівалентний універсум $\Type$, можливі також інші прямі визначення, наприклад, $K(G,1)$ еквівалентний типу малих $G$-торсорів.)
Нехай $\susp X$ позначає надбудову $X$, тобто гомотопічний пушаут $1 \leftarrow X \to 1$. Для абелевої групи $A$ тепер можна індуктивно визначити $K(A,n+1) := \trunc{\susp K(A,n)}_{n+1}$. Тоді ми маємо наступний результат~\cite{FinsterLicata2014}.

\begin{theorem}\label{thm:EM-spaces}
  Нехай $G$ — група і $n \ge 1$, і припустимо, що $G$ абелева, коли $n > 1$. Простір $K(G,n)$ є $(n-1)$-зв'язним і $n$-усіченим, і існує ізоморфізм груп $\pi_n K(G,n) \equiv G$.
\end{theorem}

У деяких наших неформальних аргументах ми використовуємо теорему про десент для пушаутів,\footnote{Згадаємо з~\cite[\S6.1.3]{LurieHTT}, слідуючи ідеям Чарльза Резка, що ми можемо \emph{визначити} $\infty$-топоси серед локально декартово замкнених $\infty$-категорій як ті, коліміти яких є \emph{ван Кампеновими}, тобто задовольняють десент.} яка стверджує, що для комутативного куба типів
\begin{equation}\label{eq:cube}
\begin{tikzcd}
& A_{11} \arrow[dl] \arrow[dr] \arrow[d] \\
A_{10} \arrow[d] & B_{11} \arrow[dl] \arrow[dr] & A_{01} \arrow[dl,crossing over] \arrow[d] \\
B_{10} \arrow[dr] & A_{00} \arrow[d] \arrow[from=ul,crossing over] & B_{01} \arrow[dl] \\
& B_{00},
\end{tikzcd}
\end{equation}
якщо нижній квадрат є пушаутом, а вертикальні квадрати — пуллбеками, то верхній квадрат також є пушаутом. Ми використаємо наступне невелике узагальнення.

\begin{theorem}
Розглянемо комутативний куб типів як у~\eqref{eq:cube}, і припустимо, що вертикальні квадрати є пуллбек-квадратами. Тоді квадрат
\begin{equation*}
\begin{tikzcd}
A_{10} \sqcup^{A_{11}} A_{01} \arrow[r] \arrow[d] & A_{00} \arrow[d] \\
B_{10} \sqcup^{B_{11}} B_{01} \arrow[r] & B_{00}
\end{tikzcd}
\end{equation*}
є пуллбек-квадратом.
\end{theorem}

\begin{proof}
Достатньо показати, що пуллбек
\begin{equation*}
(B_{10} \sqcup^{B_{11}} B_{01}) \times_{B_{00}} A_{00}
\end{equation*}
має універсальну властивість пушаута. Це випливає з теореми про десент, оскільки за лемою про вставку для пуллбеків ми також маємо, що вертикальні квадрати в кубі
\begin{equation*}
\begin{tikzcd}
& A_{11} \arrow[dl] \arrow[d] \arrow[dr] \\
A_{10} \arrow[d] & B_{11} \arrow[dl] \arrow[dr] & A_{01} \arrow[dl,crossing over] \arrow[d] \\
B_{10} \arrow[dr] & (B_{10} \sqcup^{B_{11}} B_{01}) \times_{B_{00}} A_{00} \arrow[d] \arrow[from=ul,crossing over] & B_{01} \arrow[dl] \\
& B_{10} \sqcup^{B_{11}} B_{01}
\end{tikzcd}
\end{equation*}
є пуллбек-квадратами.
\end{proof}

У формалізації аргументи з використанням десенту зручніше робити через еквівалентний принцип, зафіксований формально як лема про сплощення~\cite[\S6.12]{TheBook}.

\section{Вищі групи}
\label{sec:higher-groups}

Згадаємо, що типи в HoTT можна розглядати як $\infty$-групоїди: елементи — це об'єкти, шляхи — морфізми, вищі шляхи — вищі морфізми і т. д.

Звідси випливає, що \emph{точковий зв'язний} типи $A$ можна розглядати як вищі групи з \emph{носієм} $\Loop A := (\pt =_A \pt)$. Нейтральний елемент — це шлях тотожності, групова операція задається композицією шляхів, а вищі шляхи засвідчують закони одиниці та асоціативності. Звичайно, ці вищі шляхи самі по собі підлягають подальшим законам і т. д., але краса теоретико-типового визначення полягає в тому, що нам не потрібно про це турбуватися: всі (вищі) закони випливають з правил типів тотожності.

Пишучи $G$ для носія, зазвичай пишуть $BG$ для точкового зв'язного типу, такого що $G = \Loop BG$. Ми називаємо $BG$ \emph{розпетлюванням} $G$. Нехай ми пишемо
\begin{align*}
  \infty\Grp
  &:= (G : \Type) \times (BG : \Type_\pt^{>0}) \times (G \equiv \Loop BG) \\
  &\phantom{:}\equiv (G : \Type_\pt) \times (BG : \Type_\pt^{>0}) \times (G \equiv_\pt \Loop BG) \\
  &\phantom{:}\equiv \Type_\pt^{>0}
\end{align*}
для типу вищих груп, або \emph{$\infty$-груп}. Зауважте, що для $G : \infty\Grp$ ми також маємо $G : \Type$, використовуючи першу проекцію як коерцію. Використовуючи останнє визначення, це відображення простору петель, а не звичайна коерція!

Ми повертаємо звичайні групи на рівні множин, вимагаючи, щоб $G$ був $0$-типом, або еквівалентно, щоб $BG$ був $1$-типом. Це веде нас до введення
\begin{align*}
  n\Grp
  &:= (G : \Type_\pt^{<n}) \times (BG : \Type_\pt^{>0})
    \times (G \equiv_\pt \Loop BG) \\
  &\phantom{:}\equiv \Type_\pt^{>0,\le n}
\end{align*}
для типу \emph{групових} (групоподібних) $(n-1)$-групоїдів, також відомих як \emph{$n$-групи}. Для $G : 1\Grp$ групи на рівні множин, ми маємо $BG = K(G,1)$.

Наприклад, цілі числа $\bZ$ як адитивна група з цієї перспективи представлені їх розпетлюванням $\BZ = \bS^1$, тобто колом.

Звичайно, подвійні простори петель поводяться навіть краще, ніж просто простори петель (наприклад, вони комутативні з точністю до гомотопії за аргументом Екманна-Хілтона~\cite[Теорема~2.1.6]{TheBook}). Скажемо, що тип $G$ є \emph{$k$-кратно груповим}, якщо ми маємо $k$-кратне розпетлювання, $B^k G : \Type_\pt^{\ge k}$, таке що $G = \Loop^k B^k G$.

Змішуючи два напрямки, введемо тип
\begin{align*}
  (n,k)\GType
  &:= (G : \Type_\pt^{\le n}) \times (B^k G : \Type_\pt^{\ge k})
    \times (G \equiv_\pt \Loop^k B^k G) \\
  &\omit$\phantom{:}\equiv \Type_\pt^{\ge k,\le n+k}$\hfill\text{\leanref{GType\_equiv}}
\end{align*}
для типу \emph{$k$-кратно групових $n$-групоїдів}.\footnote{Це називається $n\Type_k$ у~\cite{BaezDolan1998}, але тут ми надаємо однакову вагу $n$ та $k$, і додаємо «G», щоб вказати на групову структуру.} (Ми дозволяємо приймати $n = \infty$, і в цьому випадку вимога усічення просто відкидається. \leanref{InfGType\_equiv})

Зауважте, що $n\Grp = (n-1,1)\GType$. Цей зсув в індексації трохи дратує, але ми зберігаємо його, щоб бути послідовними з літературою.

Оскільки існують забуваючи відображення
\begin{equation*}
(n,k+1)\GType \to (n,k)\GType
\end{equation*}
задані $B^{k+1} G \mapsto \Loop B^{k+1} G$, ми можемо також дозволити $k$ бути нескінченним, $k = \omega$, встановивши
\begin{align*}
(n,\omega)\GType & := \lim\nolimits_k (n,k)\GType \\
&\phantom{:}\equiv \bigl(B^{-} G : (k : \bN) \to \Type_\pt^{\ge k,\le n+k}\bigr) \\
  &\qquad \times \bigl((k : \bN) \to B^k G \equiv_\pt \Loop B^{k+1} G \bigr).
\end{align*}
У~\autoref{sec:stabilization} ми доводимо теорему про стабілізацію (\autoref{thm:stabilization}), з якої випливає, що $(n,\omega)\GType = (n,k)\GType$ для $k \ge n+2$.

Коли $(n,k) = (\infty,\omega)$, це тип стабільно групових $\infty$-груп, також відомих як \emph{зв'язні спектри}. Якщо ми також послабимо вимогу зв'язності, ми отримаємо тип усіх спектрів, і ми можемо думати про спектр як про свого роду $\infty$-групоїд з $k$-морфізмами для всіх $k \in \bZ$.

Класи вищих груп підсумовані в~\autoref{tab:periodic}. Ми доведемо правильність стовпця $n=0$ у~\autoref{sec:n=0}.

\section{Елементарна теорія}
\label{sec:elementary-theory}

Для \emph{будь-якого} типу об'єктів $A$, будь-який $a : A$ має \emph{групу автоморфізмів} $\Aut_A a := \Aut a := (a = a)$ з $\BAut a = \im(a : 1 \to A) = (x : A) \times \trunc{a = x}_{-1}$ (зв'язна компонента $A$ в $a$). Очевидно, якщо $A$ є $(n+1)$-усіченим, то $\BAut a$ також є таким, і тому $\Aut a$ є $n$-усіченим, а отже, $(n+1)$-групою.

Переходячи через гомотопічну гіпотезу, для кожного точкового типу $(X,x)$ ми маємо \emph{фундаментальну $\infty$-групу $X$}, $\Pi_\infty(X,x) := \Aut x$. Її $(n-1)$-усічення (приклад декатегоризації, див.~\autoref{sec:stabilization}) є \emph{фундаментальною $n$-групою $X$}, $\Pi_n(X,x)$, з відповідним розпетлюванням $\mathrm{B}\Pi_n(X,x) = \trunc{\BAut x}_n$.

Якщо ми візьмемо $A = \Set$, ми отримаємо звичайні симетричні групи $\Sym_n := \Aut(\Fin n)$, де $\Fin n$ — множина з $n$ елементами. (Зауважте, що $\BS_n = \BAut(\Fin n)$ — це тип усіх $n$-елементних множин.) Ми наведемо подальші конструкції, пов'язані зі звичайними групами, у~\autoref{sec:perspectives}.

\subsection{Гомоморфізми та спряження}
\label{sec:homomorphisms}

Гомоморфізм між вищими групами — це будь-яка функція, яку можна відповідним чином розпетлювати. Для $G,H : (n,k)\GType$ ми визначаємо
\begin{align*}
\hom_{(n,k)}(G,H) & := 
(h: G \to_\pt H)\times (B^k h: B^kG \to_\pt B^kH) \\
& \qquad \times (\Loop^k(B^k h) \sim_\pt h) \\
&\phantom{:}\equiv (B^k h: B^kG \to_\pt B^kH).
\end{align*}
Для (зв'язних) спектрів нам потрібні точкові відображення між усіма розпетлюваннями та точкові гомотопії, що показують їх когерентність.

Зауважте, що якщо $h,k : G \to H$ — гомоморфізми між групами на рівні множин, то $h$ і $k$ \emph{спряжені}, якщо $Bh, Bk : BG \to_\pt BH$ є \emph{вільно} гомотопними (тобто рівними як відображення $BG \to BH$).

Також зауважте, що $\pi_j(B^kG \to_\pt B^kH) \equiv \trunc{B^kG \to_\pt \Loop^kB^kH}_0 \equiv \trunc{\Sigma^jB^kG \to_\pt B^kH}_0 = 0$ для $j+k-1 \ge n+k$, тобто для $j > n$, тому це свідчить про те, що $\hom_{(n,k)}(G,H)$ є $n$-усіченим. (Розрахунок підтверджує це для компоненти тотожності.) Щоб довести це, нам потрібно використовувати індукцію, використовуючи визначення $n$-усіченого. Якщо $f : \hom_{(n,k)}(G,H)$, то його тип тотожності до самого себе еквівалентний $\bigl(\alpha : (z : B^kG) \to (f\,z = f\,z)\bigr) \times \bigl(\alpha\,\pt \concat g_\pt = f_\pt\bigr)$. Цей тип більше не є типом точкових відображень, а скоріше типом точкових перетинів розшарування точкових типів.

\begin{definition}
  Якщо $X : \Type_\pt$ і $Y : X \to \Type_\pt$, то ми вводимо тип \emph{точкових перетинів},
  \[
    (x : X) \to_\pt Y\,x
    := \bigl(s : (x : X) \to Y\,x\bigr)
    \times \bigl(s\,\pt = \pt\bigr).
  \]
  Цей тип сам є точковим завдяки тривіальному перетину $\lambda x, \pt$.
\end{definition}

\begin{theorem}
  Нехай $X : \Type_\pt^{\ge k}$ — $(k-1)$-зв'язний точковий тип для деякого $k \ge 0$, і нехай $Y : X \to \Type_\pt^{\le n+k}$ — розшарування $(n+k)$-усічених точкових типів для деякого $n \ge -1$. Тоді тип точкових перетинів, $(x : X) \to_\pt Y\,x$, є $n$-усіченим. \leanref{is\_trunc\_ppi\_of\_is\_conn}
\end{theorem}

\begin{proof}
  Доведення проводиться індукцією по $n$. Для базового випадку $n=-1$ ми маємо показати, що тип точкових перетинів є чистою пропозицією. Оскільки він точковий, він фактично має бути контрактним. Центром контракції є тривіальний перетин $s_0$. Якщо $s$ — інший перетин, то ми отримуємо точкову гомотопію від $s$ до $s_0$ з принципу елімінації для точкових зв'язних типів~\cite[Лема~7.5.7]{TheBook}, оскільки типи $s\,x = s_0\,x$ є $(k-2)$-усіченими.

  Щоб показати результат для $n+1$, приймаючи випадок $n$ як індукційну гіпотезу, достатньо показати для будь-якого точкового перетину $s$, що його тип тотожності до самого себе є $n$-усіченим. Але цей тип еквівалентний $(x : X) \to_\pt \Loop(Y\,x, s\,x)$, що знову є типом точкових перетинів, і тут ми можемо застосувати індукційну гіпотезу.
\end{proof}

\begin{corollary}
  Нехай $k \ge 0$ і $n \ge -1$. Якщо $X$ є $(k-1)$-зв'язним, а $Y$ — $(n+k)$-усіченим, то тип точкових відображень $X \to_\pt Y$ є $n$-усіченим. Зокрема, $\hom_{(n,k)}(G,H)$ є $n$-типом для $G,H : (n,k)\GType$.
\end{corollary}

\begin{corollary}
  Тип $(n,k)\GType$ є $(n+1)$-усіченим. \leanref{is\_\allowbreak trunc\_\allowbreak GType}
\end{corollary}

\begin{proof}
  Це випливає безпосередньо з попереднього наслідку, оскільки тип еквівалентностей $G \equiv H$ є підтипом гомоморфізмів від $G$ до $H$.
\end{proof}

Якщо $k \ge n+2$ (тобто ми знаходимося в стабільному діапазоні), то $\hom_{(n,k)}(G,H)$ стає стабільно груповим $n$-групоїдом. Це узагальнює той факт, що гомоморфізми між абелевими групами утворюють абелеву групу.

Група автоморфізмів $\Aut G$ вищої групи $G : (n,k)\GType$ знаходиться в $(n,1)\GType$. Це еквівалентно групі автоморфізмів точкового типу $B^k G$. Але ми також можемо забути базову точку і розглянути групу автоморфізмів $\Aut^c G$ типу $B^k G : \Type^{\ge k,\le n+k}$. Тепер це дозволяє (вищі) спряження. Ми визначаємо \emph{узагальнений центр} $G$ як $ZG := \Loop^k \Aut^c G : (n,k+1)\GType$ (узагальнюючи центр групи на рівні множин, див. нижче в~\autoref{sec:center}).

\subsection{Дії груп}
\label{sec:actions}

У цьому розділі ми розглядаємо фіксовану групу $G : \GType$ з розпетлюванням $BG$. \emph{Дія} $G$ на деякому об'єкті типу $A$ — це просто функція $X : BG \to A$. Об'єкт дії — це $X(\pt) : A$, і може бути зручно вважати оцінку в $\pt : BG$ коерцією від дій типу $A$ до $A$. Оснастити $a : A$ дією $G$ — це дати дію $X : BG \to A$ з $X(\pt) = a$. \emph{Тривіальна дія} — це постійна функція в $a$. Очевидно, дія $G$ на $a : A$ — це те саме, що гомоморфізм $G \to \Aut_A a$.

Якщо $A$ — універсум типів, то ми маємо дії на типах $X : BG \to \Type$. Ці $G$-типи, таким чином, є просто типами в контексті $BG$. Відображення $G$-типів від $X$ до $Y$ — це просто функція $\alpha : (z : BG) \to X(z) \to Y(z)$.

Якщо $X$ — $G$-тип, то ми можемо сформувати
\begin{description}
\item[інваріанти] $X^{hG} := (z : BG) \to X(z)$, також відомі як \emph{гомотопічні нерухомі точки}, та
\item[коінваріанти] $X_{hG} := (z : BG) \times X(z)$, які також відомі як \emph{гомотопічний простір орбіт} або \emph{гомотопічна частка} $X \dblslash G$.
\end{description}
Легко бачити, що ці конструкції є відповідно правим і лівим спряженими до функтора, який надсилає тип $X$ до тривіальної дії $G$ на $X$, $X^{\mathrm{triv}} : BG \to \Type$, яка є просто постійним сімейством в $X$. Дійсно, спряження — це просто звичайні аргументи перестановки та еквівалентності (роз)карріювання, для $Y : \Type$,
\begin{align*}
  \hom(Y, X^{hG})
  &= X \to (z : BG) \to Y(z)
    \equiv (z : BG) \to X \to Y(z) \\
  &\equiv \hom(X^{\mathrm{triv}}, Y), \\
  \hom(X_{hG}, Y)
  &= \big((z : BG) \times X(z)\bigr) \to Y
    \equiv (z : BG) \to X(z) \to Y \\
  &\equiv \hom(X, Y^{\mathrm{triv}}).
\end{align*}
Якщо ми думаємо про дію $X : BG \to \Type$ як про діаграму зі значеннями в типах на $BG$, це означає, що гомотопічні нерухомі точки та гомотопічний простір орбіт утворюють гомотопічну границю та гомотопічну кограницю цієї діаграми відповідно.

\begin{proposition}
  \label{prop:action-hom-pullback}
  Нехай $f : H \to G$ — гомоморфізм вищих груп з розпетлюванням $Bf : BH \to_\pt BG$, і нехай $\alpha : \hom(X,Y)$ — відображення $G$-типів. Компонуючи з $f$, ми також можемо розглядати $X$ і $Y$ як $H$-типи, і в цьому випадку ми отримуємо гомотопічний пуллбек-квадрат:
  \[
    \begin{tikzcd}
      X_{hH} \ar[r]\ar[d] & Y_{hH}\ar[d] \\
      X_{hG} \ar[r]       & Y_{hG}.
    \end{tikzcd}
  \]
\end{proposition}

\begin{proof}
  Вертикальні відображення індукуються $Bf$, а горизонтальні — $\alpha$. Тип гомотопічного пуллбека $C$ обчислюється як
  \begin{align*}
    C
    &\equiv (z : BG) \times (x : X\,z)
      \times (w : BH) \times (y : Y(Bf\,w)) \\
    &\qquad \times (z = Bf\,w) \times (y = \alpha\,z\,x) \\
    &\equiv (w : BH) \times (x : X(Bf\,w))
      = X_{hH},
  \end{align*}
  і за цієї еквівалентності верхнє та ліве відображення є канонічними.
\end{proof}

Кожна група $G$ несе дві канонічні дії на собі:
\begin{description}
\item[права дія] $G : BG \to \Type$, $G(x) = (\pt = x)$, та
\item[приєднана дія] $G^{\mathrm{ad}} : BG \to \Type$, $G^{\mathrm{ad}}(x) = (x = x)$ (шляхом спряження).
\end{description}

Ми маємо $1 \dblslash G = BG$, $G \dblslash G = 1$ і $G^{\mathrm{ad}} \dblslash G = LBG := (\bS^1 \to BG)$, вільний простір петель $BG$. Пам’ятаючи, що $\BZ = \bS^1$, ми бачимо, що $G^{\mathrm{ad}} = (\BZ \to BG)$, тобто класи спряженості гомоморфізмів з $\bZ$ в $G$. Оскільки цілі числа є вільною (вищою) групою на одному генераторі, це просто класи спряженості елементів $G$. Але це саме те, що ми повинні отримати для гомотопічних орбіт $G$ під дією спряження.

Попередня пропозиція має цікавий наслідок:
\begin{corollary}
  \label{cor:hfiber-hom}
  Якщо $f : H \to G$ — гомоморфізм вищих груп, то $G \dblslash H$ еквівалентно гомотопічному волокну розпетлювання $Bf : BH \to_\pt BG$, де $H$ діє на $G$ через $f$-індуковану праву дію.
\end{corollary}

\begin{proof}
  Ми застосовуємо Пропозицію~\ref{prop:action-hom-pullback} з $\alpha : G \to 1$, що є канонічним відображенням від правої дії $G$ до дії $G$ на одиничному типі. Тоді квадрат стає:
  \[
    \begin{tikzcd}[baseline=(O.base)]
      G \dblslash H \ar[r]\ar[d] & BH\ar[d] \\
      |[alias=O]| 1 \ar[r]       & BG
    \end{tikzcd}\qedhere
  \]
\end{proof}

За визначенням, $BG$ класифікує \emph{головні $G$-розшарування}: пуллбеки правої дії $G$. Тобто головне $G$-розшарування над типом $A$ — це сімейство типів $F : X \to \Type$, представлене відображенням $\chi : A \to BG$, таким що $F(x) \equiv (\pt = \chi(x))$ для всіх $x : X$.

Наприклад, для кожної вищої групи $G$ ми маємо відповідне розшарування Хопфа $\susp G \to \Type$, представлене відображенням $\chi_H : \susp G \to BG$, що відповідає за спряженістю петля-надбудова тотожному відображенню на $G$. (Це конкретне розшарування можна визначити, використовуючи лише індуковану структуру $H$-простору на $G$.)

Ця перспектива лежить в основі побудови першим та третім іменованими авторами дійсних проективних просторів у гомотопічній теорії типів~\cite{realprojective}. Послідовності волокон $\bS^0 \to \bS^n \to \RP^n$ є головними розшаруваннями для групи з 2-х елементів $\bS^0 = \Sym_2$ з розпетлюванням $\BS_2 \equiv \RP^\infty$, типом 2-елементних типів.

\subsection{Назад до центру}
\label{sec:center}

Ми згадували узагальнений центр вище і стверджували, що він узагальнює звичайне поняття центру групи. Дійсно, якщо $G : 1\Grp$ — група на рівні множин, то елемент $ZG$ відповідає елементу $\Loop^2 \BAut^c G$, або еквівалентно, відображенню з 2-сфери $\bS^2$ в $\Type$, що надсилає базову точку в $BG$. За універсальною властивістю $\bS^2$ як ВІТ, це знову відповідає гомотопії від тотожності на $BG$ до самої себе, $c : (z : BG) \to z = z$. Це саме гомотопічна нерухома точка приєднаної дії $G$ на собі, тобто центральний елемент.

\subsection{Еквіваріантна гомотопічна теорія}
\label{sec:equivariant}

Фіксуємо групу $G : \GType$. Припустимо, що $G$ — це насправді (гомотопічний) тип топологічної групи. Розглянемо тип $BG \to \Type$ (малих) \emph{типів з дією $G$}. Наївно можна було б подумати, що це представляє $G$-еквіваріантні гомотопічні типи, тобто достатньо хороші\footnote{Достатньо хороші означає $G$-CW-простори. Та сама гомотопічна категорія виникає, якщо взяти всі простори з дією $G$, але тоді слабкі еквівалентності — це $G$-відображення $f : X \to Y$, які індукують слабкі еквівалентності на просторах $H$-нерухомих точок $f^H : X^H \to Y^H$ для всіх замкнутих підгруп $H$ в $G$.} топологічні простори з дією $G$, що розглядаються з точністю до $G$-еквіваріантної гомотопічної еквівалентності. Але це не так.

За теоремою Елмендорфа~\cite{Elmendorf1983}, ця гомотопічна теорія скоріше є теорією препучків (звичайних) гомотопічних типів на \emph{категорії орбіт} $\mathcal{O}_G$ групи $G$. Це повна підкатегорія категорії $G$-просторів, натягнута на однорідні простори $G/H$, де $H$ пробігає замкнуті підгрупи $G$.

Всередині категорії орбіт ми знаходимо копію групи $G$, а саме як ендоморфізми об'єкта $G/1$, що відповідає тривіальній підгрупі $1$. Отже, $G$-еквіваріантний гомотопічний тип дає тип з дією $G$ шляхом обмеження вздовж включення $BG \hookrightarrow \mathcal{O}_G$. (Тут ми розглядаємо $BG$ як (точковий і зв'язний) топологічний групоїд на одному об'єкті.)

Як зауважив Шульман~\cite{ShulmanEI}, коли $G$ — компактна група Лі, то $\mathcal{O}_G$ є інверсною EI $\infty$-категорією, і тому ми знаємо, як моделювати теорію типів у $\infty$-топосі препучків над $\mathcal{O}_G$. І в деяких простих випадках ми можемо навіть визначити цю модель внутрішньо. Наприклад, якщо $G = \bZ/p\bZ$ — циклічна група простого порядку, то малий $G$-еквіваріантний тип складається з типу з дією $G$, $X : BG \to \Type$, разом з іншим сімейством типів $X^G : X^{hG} \to \Type$, де $X^G$ дає для кожної гомотопічної нерухомої точки тип доказів або «спеціальних причин», чому цю точку слід вважати нерухомою~\cite[7.6]{ShulmanEI}. Отже, повний простір $X^G$ — це тип фактичних нерухомих точок, а проекція на $X^{hG}$ реалізує відображення від фактичних нерухомих точок до гомотопічних нерухомих точок.

Навіть не звертаючись до категорії орбіт, ми можемо дещо сказати про топологічні групи через їхні класифікуючі типи в теорії типів. Наприклад~\cite{Camarena}, якщо $f : H \to G$ ін'єктивне, то гомотопічне волокно $Bf$ за Наслідком~\ref{cor:hfiber-hom} є гомотопічним простором орбіт $G \dblslash H$, який у цьому випадку є просто простором суміжних класів $G/H$, і, отже, в теорії типів представляє гомотопічний тип цього простору суміжних класів. І якщо
\[
  1 \to K \to G \to H \to 1
\]
є короткою точною послідовністю топологічних груп, то $BK \to BG \to BH$ є послідовністю розшарування, тобто ми можемо відновити розпетлювання $BK$ групи $K$ як гомотопічне волокно відображення $BG \to BH$.

\subsection{Деякі елементарні конструкції}
\label{sec:elementary-constructions}

Якщо нам дано гомоморфізм $\varphi : H \to \Aut(N)$, представлений точковим відображенням $B\varphi : BH \to_\pt \BAut_\pt(BN)$, де $\BAut_\pt(BN)$ — тип точкових типів, лише еквівалентних $BN$, ми можемо побудувати нову групу, \emph{напівпрямий добуток}, $G := H \ltimes_\varphi N$ з класифікуючим типом $BG := (z : BH) \times (B\varphi\,z)$. Тип $BG$ справді точковий (за парою базової точки $\pt$ в $BG$ і базової точки в точковому типі $B\varphi(\pt)$), і зв'язний, а отже, представляє вищу групу $G$. Елемент $g$ задається парою елемента $h : H$ та ідентифікацією $g \concat \pt = \pt$ в $B\varphi(\pt) \equiv_\pt BN$. Але оскільки дія здійснюється через точкові відображення, другий компонент еквівалентний ідентифікації $\pt = \pt$ в $BN$, тобто елементу $N$. За цією еквівалентністю добуток $(h,n)$ і $(h',n')$ дійсно є $(h \concat h', n \concat \varphi(h)(n'))$.

Як окремий випадок ми отримуємо \emph{прямий добуток}, коли $\varphi$ є тривіальною дією. Тут $B(H \times N) \equiv BH \times BN$.

Як інший окремий випадок ми отримуємо \emph{сплетіння} $N \wr \Sym_n$ групи $N$ і симетричної групи $\Sym_n$. Тут $\Sym_n$ діє на прямий степінь $N^{\Fin n}$, переставляючи фактори. Дійсно, використовуючи представлення $\BS_n$ як типу $n$-елементних типів, відображення $B\varphi$ — це просто $A \mapsto (A \to BN)$. Отже, розпетлювання сплетіння $G := N \wr \Sym_n$ — це просто $BG := (A : BS_n) \times (A \to BN)$.

\section{Групи на рівні множин}
\label{sec:n=0}
У цьому розділі ми наводимо доказ того, що стовпець $n=0$ у~\autoref{tab:periodic} правильний. Зауважте, що для $n=0$ типи гомоморфізмів $\hom_{(0,k)}(G,H)$ є множинами, що означає, що $(0,k)\GType$ утворює 1-категорію. Нехай $\Group$ — категорія звичайних груп на рівні множин (множина з множенням, оберненим елементом та одиницею, що задовольняють закони групи), а $\AbGroup$ — категорія абелевих груп.

\begin{theorem}
  Ми маємо наступні еквівалентності категорій {\normalfont(}для $k \ge 2${\normalfont):}
  \begin{align*}
    (0,1)\GType &\equiv \Group;&\text{\leanref{cGType\_equivalence\_Grp}}&\\
    (0,k)\GType &\equiv \AbGroup.&\text{\leanref{cGType\_equivalence\_AbGrp}}&
  \end{align*}
\end{theorem}

Оскільки ця теорема була формалізована, ми не будемо наводити всі деталі доведення.

\begin{proof}
  Нехай $k \ge 1$ і $G$ — група, яка є абелевою, якщо $k > 1$, і нехай $X : \Type_\pt^{\ge k, \le k}$. Якщо ми маємо груповий гомоморфізм $\phi : G \to \Loop^k X$, ми отримуємо відображення $e_\phi^k : K(G,k) \to_\pt X$. Для $k=1$ це випливає безпосередньо з принципу індукції $K(G,1)$. Для $k > 1$ ми можемо визначити груповий гомоморфізм $\widetilde\phi$ як композицію $G \xrightarrow{\phi} \Loop^k X \equiv \Loop^{k-1}(\Loop X)$ і застосувати індукційну гіпотезу, щоб отримати відображення $e_{\widetilde\phi}^{k-1} : K(G,k-1) \to_\pt \Loop X$. За спряженістю $\Sigma \dashv \Loop$ ми отримуємо точкове відображення $\Sigma K(G,k-1) \to_\pt X$, а за принципом елімінації усічення ми отримуємо відображення $K(G,k) = \trunc{\Sigma K(G,k-1)}_k \to_\pt X$.

  Тепер ми можемо показати, що $\Loop^k e_\phi^k$ є очікуваним відображенням, тобто наступна діаграма комутує, але ми опускаємо це доведення тут.
  \begin{center}
    \begin{tikzpicture}[auto,node distance=2cm,
        thick,main node/.style={font=\sffamily\bfseries},text height=1.5ex]
      \node[main node] (OK) at (0,0) {$\Loop^n K(G,k)$};
      \node[main node] (G)  at (3,0) {$G$};
      \node[main node] (OX) at (1.5,-1.5) {$\Loop^n X$};
           \path[every node/.style={font=\sffamily\small}]
           (OK) edge [->] node {$\sim$} (G)
                edge [->] node [below left] {$\Loop^n e_\phi^k$} (OX)
           (OX) edge [->] node [below right] {$\phi$} (G);
    \end{tikzpicture}
  \end{center}
  Тепер, якщо $\phi$ — ізоморфізм груп, за теоремою Вайтхеда для усічених типів~\cite[Теорема 8.8.3]{TheBook} ми знаємо, що $e_\phi^k$ є еквівалентністю, оскільки воно індукує еквівалентність на всіх гомотопічних групах (тривіально на рівнях, відмінних від $k$). Ми також можемо показати, що $e_\phi^k$ натуральне за $\phi$.

  Зауважте, що якщо ми маємо груповий гомоморфізм $\psi : G \to G'$, ми також отримуємо груповий гомоморфізм $G \to \Loop^k K(G',k)$, і за допомогою вищеописаної конструкції ми отримуємо точкове відображення $K(\psi,k) : K(G,k) \to_\pt K(G',k)$. Це функторіально, що випливає з натуральності $e_\phi^k$.

  Нарешті, ми можемо побудувати еквівалентність явно. Ми маємо функтор $\pi_k : (0,k)\GType \to \AbGroup$, який надсилає $G$ в $\pi_k BG$. Навпаки, ми маємо функтор $K({-},k) : \AbGroup \to (0,k)\GType$. Ми маємо натуральні ізоморфізми $\pi_k K(G,k) \equiv G$ за~\autoref{thm:EM-spaces} і $K(\pi_k X,k) \equiv_\pt X$ за застосуванням Вайтхеда, описаного вище. Конструкція точно така ж для $k=1$ після заміни $\AbGroup$ на $\Group$.
\end{proof}

\section{Стабілізація}
\label{sec:stabilization}

У цьому розділі ми обговорюємо деякі конструкції з вищими групами~\cite{BaezDolan1998}. Ми наведемо дії на носіях та розпетлюваннях, але опустимо третій компонент, точкову еквівалентність, для зручності читання. Ми рекомендуємо тримати~\autoref{tab:periodic} в голові під час цих конструкцій.
\begin{description}
\item[декатегоризація] $\Decat : (n,k)\GType \to (n-1,k)\GType$\\
  $\angled{G,B^kG} \mapsto \angled{\trunc G_{n-1}, \trunc{B^kG}_{n+k-1}}$
\item[дискретна категоризація] $\Disc : (n,k)\GType \to (n+1,k)\GType$ \\
  $\angled{G,B^kG} \mapsto \angled{G,B^kG}$
\end{description}
Ці функтори роблять $(n,k)\GType$ рефлексивною під-$(\infty,1)$-категорією $(n+1,k)\GType$. Тобто існує спряженість ${\Decat} \dashv {\Disc}$ \leanref{Decat\_adjoint\_Disc}\footnote{У формалізації натуральність спряженості є окремим твердженням, \leanref{Decat\_adjoint\_Disc\_natural}. Це також вірно для інших спряженостей.} таке, що коодиниця індукує ізоморфізм ${\Decat} \circ {\Disc} = \id$ \leanref{Decat\_Disc}. Ці властивості є прямими наслідками універсальної властивості усічення.

Існують також ітеровані версії цих функторів.
\begin{description}
  \item[$\infty$-декатегоризація] $\iDecat : (\infty,k)\GType \to (n,k)\GType$\\
    $\angled{G,B^kG} \mapsto \angled{\trunc G_{n}, \trunc{B^kG}_{n+k}}$
  \item[дискретна $\infty$-категоризація] $\iDisc : (n,k)\GType \to (\infty,k)\GType$ \\
    $\angled{G,B^kG} \mapsto \angled{G,B^kG}$
\end{description}
Ці функтори мають ті ж властивості: ${\iDecat} \dashv {\iDisc}$ \leanref{InfDecat\_adjoint\_InfDisc} таке, що коодиниця індукує ізоморфізм ${\iDecat} \circ {\iDisc} = \id$ \leanref{InfDecat\_InfDisc}.

Для наступних конструкцій нам потрібні такі властивості.
\begin{definition}
  Для $A : \Type_\pt$ ми визначаємо \emph{$n$-зв'язне накриття} $A$ як $A{\angled n} := \fib(A \to \trunc A_n)$. Ми маємо проекцію $p_1 : A{\angled n} \to_\pt A$.
\end{definition}
\begin{lemma} \label{lem:connected-cover-univ}
  Універсальна властивість $n$-зв'язного накриття стверджує наступне. Для будь-якого $n$-зв'язного точкового типу $B$ точкове відображення
  $$(B \to_\pt A{\angled n}) \to_\pt (B \to_\pt A),$$
  задане постпозицією з $p_1$, є еквівалентністю. \leanref{connect\_intro\_pequiv}
\end{lemma}
\begin{proof}
  Дано відображення $f : B \to_\pt A$, ми можемо сформувати відображення $\widetilde f : B \to A{\angled n}$. Спочатку зауважте, що для $b : B$ тип $\tr{fb}_n =_{\trunc A_n} \tr{\pt}_n$ є $(n-1)$-усіченим і заселеним для $b = \pt$. Оскільки $B$ є $n$-зв'язним, універсальна властивість для зв'язних типів показує, що ми можемо побудувати $qb : \tr{fb}_n = \tr{\pt}_n$ для всіх $b$ таких, що $q_0 : qb_0 \concat \ap_{\tr{\blank}_n}(f_0) = 1$. Тоді ми можемо визначити відображення $\widetilde f(b) := (fb, qb)$. Тепер $\widetilde f$ точкове, тому що $(f_0,q_0) : (fb_0,qb_0) = (a_0,1)$.

  Тепер ми показуємо, що це дійсно зворотне до даного відображення. З одного боку, нам потрібно показати, що якщо $f : B \to_\pt A$, то $p_1 \circ \widetilde f = f$. Базові функції рівні, оскільки вони обидві надсилають $b$ в $f(b)$. Вони поважають точки однаково, тому що $\ap_{p_1}(\widetilde f_0) = f_0$. Доказ того, що інша композиція є тотожністю, випливає з обчислення з використанням волокон і зв'язності, яке ми тут опускаємо, але його можна знайти у формалізації.
\end{proof}
Наступна рефлексивна під-$(\infty,1)$-категорія утворюється шляхом запетлювання та розпетлювання.
\begin{description}
\item[запетлювання] $\Loop : (n,k)\GType \to (n-1,k+1)\GType$ \\
  $\angled{G,B^kG} \mapsto \angled{\Loop G,B^kG{\angled k}}$
\item[розпетлювання] $\B : (n,k)\GType \to (n+1,k-1)\GType$ \\
  $\angled{G,B^kG} \mapsto \angled{\Loop^{k-1}B^kG,B^kG}$
\end{description}
Ми маємо ${\B} \dashv {\Loop}$ \leanref{Deloop\_adjoint\_Loop}, що випливає з Леми~\ref{lem:connected-cover-univ} та $\Loop \circ {\B} = \id$ \leanref{Loop\_Deloop}, що випливає з того факту, що $A{\angled n} = A$, якщо $A$ є $n$-зв'язним.

Остання спряжена пара функторів задається стабілізацією та забуванням. Це не утворює рефлексивну під-$(\infty,1)$-категорію.
\begin{description}
\item[забування] $F : (n,k)\GType \to (n,k-1)\GType$ \\
  $\angled{G,B^kG} \mapsto \angled{G,\Loop B^kG}$
\item[стабілізація] $S : (n,k)\GType \to (n,k+1)\GType$ \\
  $\angled{G,B^kG} \mapsto \angled{SG,\trunc{\susp B^kG}_{n+k+1}}$,\\
  де $SG = \trunc{\Loop^{k+1}\susp B^kG}_n$
\end{description}
Ми маємо спряженість ${S} \dashv {F}$ \leanref{Stabilize\_adjoint\_Forget}, яка випливає з спряженості надбудова-петля $\Sigma \dashv \Loop$ на точкових типах.

Наступною головною метою в цьому розділі є теорема про стабілізацію, яка стверджує, що знаки повтору в~\autoref{tab:periodic} виправдані.

Наступний наслідок — це майже~\cite[Лема~8.6.2]{TheBook}, але довести це в HoTT книги трохи складно. Див. формалізацію для деталей.
\begin{lemma}[Зв'язність клина]
  \label{lem:wedge-connectivity}
  Якщо $A : \Type_\pt$ є $n$-зв'язним, а $B : \Type_\pt$ є $m$-зв'язним, то відображення $A \vee B \to A \times B$ є $(n+m)$-зв'язним. \leanref{is\_conn\_fun\_prod\_of\_wedge}
\end{lemma}

Зазначимо, що існує альтернативний спосіб довести лему про зв'язність клина: згадаємо, що якщо $A$ є $n$-зв'язним, а $B$ є $m$-зв'язним, то $A \ast B$ є $(n+m+2)$-зв'язним~\cite[Теорема~6.8]{joinconstruction}. Отже, лема про зв'язність клина також є прямим наслідком наступної леми.

\begin{lemma}
Нехай $A$ і $B$ — точкові типи. Волокно включення клина $A \vee B \to A \times B$ еквівалентне $\Loop{A} \ast \Loop{B}$. 
\end{lemma}

\begin{proof}
Зауважте, що волокно $A \to A \times B$ — це $\Loop B$, волокно $B \to A \times B$ — це $\Loop A$, і, звісно, волокно $1 \to A \times B$ — це $\Loop A \times \Loop B$. Ми отримуємо комутативний куб
\begin{equation*}
\begin{tikzcd}
& \Loop A \times \Loop B \arrow[dl] \arrow[d] \arrow[dr] \\
\Loop B \arrow[d] & 1 \arrow[dl] \arrow[dr] & \Loop A \arrow[dl,crossing over] \arrow[d] \\
A \arrow[dr] & 1 \arrow[d] \arrow[from=ul,crossing over] & B \arrow[dl] \\
& A \times B
\end{tikzcd}
\end{equation*}
в якому вертикальні квадрати є пуллбек-квадратами. 

За теоремою про десент для пушаутів тепер випливає, що $\Loop A \ast \Loop B$ є волокном включення клина.
\end{proof}

Другим основним інструментом, необхідним для теореми про стабілізацію, є:
\begin{theorem}[Фройденталя]
  Якщо $A : \Type_\pt^{>n}$ з $n \ge 0$, то відображення $A \to \Loop\susp A$ є $2n$-зв'язним.
\end{theorem}
Це~\cite[Теорема~8.6.4]{TheBook}.

Останній будівельний блок, який нам потрібен:
\begin{lemma}
  Існує пуллбек-квадрат
  \[
    \begin{tikzcd}
      \susp\Loop A \ar[d,"\varepsilon_A"']\ar[r] & A \vee A \ar[d] \\
      A \ar[r,"\Delta"'] & A \times A
    \end{tikzcd}
  \]
  для будь-якого $A : \Type_\pt$.
\end{lemma}

\begin{proof}
Зауважте, що пуллбек $\Delta : A \to A \times A$ вздовж будь-якого включення $A \to A \times A$ є контрактним. Отже, ми маємо куб
\begin{equation*}
\begin{tikzcd}
& \Loop A \arrow[dl] \arrow[d] \arrow[dr] \\
1 \arrow[d] & 1 \arrow[dl] \arrow[dr] & 1 \arrow[dl,crossing over] \arrow[d] \\
A \arrow[dr] & A \arrow[d,"\Delta"] \arrow[from=ul,crossing over] & A \arrow[dl] \\
& A \times A
\end{tikzcd}
\end{equation*}
в якому всі вертикальні квадрати є пуллбек-квадратами. Тому, якщо ми зробимо пуллбек вздовж включення клина, ми отримаємо за теоремою про десент для пушаутів, що квадрат у твердженні дійсно є пуллбек-квадратом.
\end{proof}

\begin{theorem}[Стабілізація]
  \label{thm:stabilization}
  Якщо $k \ge n+2$, то $S : (n,k)\GType \to (n,k+1)\GType$ є еквівалентністю, і будь-яка $G : (n,k)\GType$ є простором нескінченних петель. \leanref{stabilization}
\end{theorem}

\begin{proof}
  Ми показуємо, що $F \circ S = \id = S \circ F : (n,k)\GType \to (n,k)\GType$ щоразу, коли $k \ge n+2$. Для першого, відображення одиниці спряженості розкладається як
  \[
    B^kG \to \Loop\susp B^kG \to \Loop\trunc{\susp B^kG}_{n+k+1}
  \]
  де перше відображення є $(2k-2)$-зв'язним за Фройденталем, а друге відображення є $(n+k)$-зв'язним. Оскільки область є $(n+k)$-усіченою, композиція є еквівалентністю кожного разу, коли $2k-2 \ge n+k$.

  Для другого, відображення коодиниці спряженості розкладається як
  \[
    \trunc{\susp\Loop B^kG}_{n+k} \to \trunc{B^kG}_{n+k} \to B^kG,
  \]
  де друге відображення є еквівалентністю. Згідно з двома лемами вище, перше відображення є $(2k-2)$-зв'язним.
\end{proof}

Наприклад, для $G : (0,2)\GType$ абелевої групи ми маємо $B^nG = K(G,n)$, простір Ейленберга-Маклейна.

Спряженість ${S} \dashv {F}$ передбачає, що вільна група на точковому наборі $X$ — це $\Loop\trunc{\susp X}_1 = \pi_1(\susp X)$. Якщо $X$ має вирішувану рівність, то $\susp X$ вже $1$-усічений. Залишається відкритою проблема, чи вірно це в загальному випадку.

Крім того, абелізація групи на рівні множин $G : 1\Grp$ — це $\pi_2(\susp BG)$. Якщо $G : (n,k)\GType$ знаходиться в стабільному діапазоні ($k \ge n+2$), то $SFG = G$.

\section{Перспективи звичайної теорії груп}
\label{sec:perspectives}

У цьому розділі ми вкажемо, як теорія вищих груп може дати новий погляд навіть на звичайну теорію груп.

З симетричних груп $\Sym_n$ ми можемо отримати інші скінченні групи, використовуючи конструкції з~\autoref{sec:elementary-constructions}. Інші групи можна побудувати безпосередньо. Наприклад, $BA_n$, класифікуючий тип знакозмінної групи, можна прийняти як тип $n$-елементних множин $X$, оснащених \emph{знаковим впорядкуванням}: це клас еквівалентності впорядкування $\Fin n \equiv X$ за модулем парних перестановок. Дійсно, існує лише два можливих знакових впорядкування, тому це визначення відповідає спочатку розгляду короткої точної послідовності
\[
  1 \to A_n \to \Sym_n \xrightarrow{\mathrm{sgn}}{} \Sym_2 \to 1
\]
де останнє відображення — це відображення знака, потім реалізації відображення знака як заданого відображенням $\mathrm{Bsgn} : \BS_n \to \BS_2$, яке приймає $n$-елементну множину в її набір знакових впорядкувань, і, нарешті, припущенню, що $BA_n$ — гомотопічне волокно $\mathrm{Bsgn}$.

Аналогічно, $BC_n$, класифікуючий тип циклічної групи з $n$ елементів, можна прийняти як тип $n$-елементних множин $X$, оснащених \emph{циклічним впорядкуванням}: клас еквівалентності впорядкування $\Fin n \equiv X$ за модулем циклічних перестановок. Але на відміну від вищенаведеного, де ми мали збіг, що $\Aut(\Sym_2) \equiv \Sym_2$, це не відповідає короткій точній послідовності. Скоріше, це відповідає послідовності
\[
  1 \to C_n \to \Sym_n \to \Aut(\Fin(n-1)) \equiv \Sym_{(n-1)!}
\]
де розпетлювання останнього відображення — це відображення з $\BS_n$ в $\BS_{(n-1)!}$, яке відображає $n$-елементну множину в набір циклічних впорядкувань, яких є $(n-1)!$ — оскільки як тільки ми зафіксуємо позицію в впорядкуванні конкретного елемента, ми вільні переставляти решту.

Як інший приклад, розглянемо відображення $p : \BS_4 \to_\pt \BS_3$, яке відображає 4-елементну множину $X$ у її набір розбиття 2-на-2, яких є 3. Використовуючи цю конструкцію, ми можемо реалізувати деякі відомі тотожності напівпрямого та сплетіння, такі як $A_4 \equiv S_2^2 \rtimes A_3$, $S_4 \equiv S_2^2 \rtimes S_3$, і, для октаедричної групи, $O_h \equiv S_2^3 \rtimes S_3 \equiv S_2 \wr S_3$.

\smallskip

Звернемося до іншого способу отримання нових груп зі старих, а саме за допомогою теорії накриттів.

\subsection{\texorpdfstring{$1$}{1}-групи та накриття}
\label{sec:covering}

Зв'язки між накриттями точкового зв'язного типу $X$ та множинами з дією фундаментальної групи $X$ вже встановлені в гомотопічній теорії типів~\cite{FavoniaHarper2016}. Згадаємо цей зв'язок і трохи розширимо його.

Для нас точковий зв'язний тип $X$ еквівалентний $\infty$-групі $G : \infty\Grp$ з розпетлюванням $BG := X$. Накриття над $BG$ — це просто сімейство типів $C : BG \to \Set$, яке потрапляє в універсум множин. Отже, згідно з нашим обговоренням дій у~\autoref{sec:actions}, це саме множина з дією $G$. Оскільки $\Set$ є 1-типом, $C$ однозначно поширюється на сімейство типів $C' : \trunc{BG}_1 \to \Set$, але $\trunc{BG}_1$ є розпетлюванням фундаментальної групи $X$, і тому $C'$ є однозначно визначеним вибором множини з дією фундаментальної групи.

Універсальне накриття — це простозв'язне накриття $BG$,
\[
  \widetilde{BG} : BG \to \Set, \quad
  z \mapsto \trunc{\pt = z}_0.
\]
Зауважте, що повний простір $\widetilde{BG}$ справді є $1$-зв'язним накриттям $BG\angled1$, оскільки $\trunc{\pt =_{BG} \pt}_0 \equiv (\tr{\pt} =_{\trunc{BG}_1} \tr{\pt})$. Також зауважте, що якщо $G$ вже є 1-групою, то це просто права дія $G$ на собі, а загалом це права дія $G$ на фундаментальній групі (тобто декатегоризація $G$) через гомоморфізм усічення від $G$ до $\pi_1(BG)$, де ми також можемо розглядати $\pi_1(BG)$ як декатегоризацію $G$ в $1\Grp$.

Загалом існує відповідність Галуа між зв'язними накриттями $BG$ та класом спряженості підгруп фундаментальної групи. Дійсно, якщо $C : BG \to \Set$ має зв'язний повний простір, то простір $(g : \trunc{BG}_1) \times C'(g)$ сам по собі є зв'язним 1-усіченим типом, і проекція на $\trunc{BG}_1$ індукує включення фундаментальних груп після того, як вибрано точку $\pt : C'(\pt)$.

\begin{theorem}[Основна теорема теорії Галуа для накриттів]
  \label{thm:galois-zero}
  $\phantom{42}$
  \begin{enumerate}
  \item Група автоморфізмів універсального накриття $\widetilde{BG}$ ізоморфна декатегоризації $G$ в 1-групу,
    \[
      \Aut(\widetilde{BG}) \equiv \Decat_1(G) \equiv \pi_1(BG).
    \]
  \item Крім того, існує контрваріантна відповідність між класами спряженості підгруп $\Decat_1(G)$ та зв'язними накриттями $BG$.
  \item Це піднімається до відповідності Галуа між підгрупами $\Decat_1(G)$ та точковими зв'язними накриттями $BG$. Нормальні підгрупи відповідають накриттям Галуа.
  \end{enumerate}
\end{theorem}

Зауважте, що універсальне накриття та тривіальне накриття (постійне в одиничному типі) є канонічно точковими, що відображає той факт, що дві тривіальні підгрупи є нормальними.

Перша частина основної теореми має чітке узагальнення на вищі групи:

\begin{theorem}[Основна теорема теорії Галуа для $n$-накриттів, частина перша]
  Група автоморфізмів універсального $n$-типового накриття $U_n(BG)$,
  \[
    U_n(BG) : BG \to \Type^{\le n},
    \quad
    z \mapsto \trunc{\pt = z}_n
  \]
  типу $BG$ ізоморфна декатегоризації $G$ в $(n+1)$-групу,
  \[
    \Aut(U_n(BG)) \equiv \Decat_{n+1}(G) \equiv \Pi_{n+1}(BG).
  \]
\end{theorem}

\begin{proof}
  Зауважте, що $\BAut(U_n(BG))$ — це образ відображення $1 \to (BG \to \Type^{\le n})$, яке надсилає канонічний елемент у $U_n(BG)$. Оскільки $BG$ зв'язний, цей образ є саме $\trunc{BG}_{n+1}$ за~\cite[Теорема~7.1]{joinconstruction}. Тоді ми закінчили, оскільки $\B\Pi_{n+1}(BG) \equiv \trunc{BG}_{n+1}$ за визначенням.
\end{proof}

Можна використовувати інші частини~\autoref{thm:galois-zero}, щоб \emph{визначити} поняття підгрупи та нормальної підгрупи для $n$-груп, які потім стають \emph{структурою на}, а не \emph{властивістю} гомоморфізму $f : K \to G$. Явно, структура \emph{нормальної підгрупи} на такому $f$ — це розпетлювання $B(G \dblslash K)$ типу $G \dblslash K$ разом з відображенням $Bq : BG \to_\pt B(G \dblslash K)$, що призводить до послідовності волокон
\begin{equation}\label{eq:normal-fiber-sequence}
  G \dblslash K \to BK \xrightarrow{Bf}{} BG \xrightarrow{Bq}{} B(G \dblslash K).
\end{equation}

\subsection{Центральні розширення та когомології груп}
\label{sec:group-cohomology}

Когомологія вищої групи $G$ — це просто когомологія її розпетлювання $BG$. Дійсно, для будь-якого спектру $A$ ми визначаємо
\[
  H_{\mathrm{Grp}}^k(G, A) := \trunc{BG \to_\pt B^kA}_0.
\]
Звичайно, щоб визначити $k$-ту групу когомологій, нам потрібне лише $k$-кратне розпетлювання $B^kA$.

Якщо $A : (\infty,2)\GType$ — плетена $\infty$-група, то ми маємо другу групу когомологій $H_{\mathrm{Grp}}^2(G, A)$, і елемент $c : BG \to_\pt B^2A$ дає початок \emph{центральному розширенню}
\[
  BA \to BH \to BG \xrightarrow{c}{} B^2A,
\]
де $BH$ — гомотопічне волокно $c$. Це піднімає у світ вищих груп звичайний результат про те, що класи ізоморфізму центральних розширень 1-групи $G$ абелевою 1-групою $A$ задаються класами когомологій в $H_{\mathrm{Grp}}^2(G, A)$.

\smallskip

У репозиторії Spectral є повна формалізація спектральної послідовності Серра для когомологій~\cite{SerreSpectralSequence}. Якщо ми маємо будь-яку послідовність волокон нормальної підгрупи для $\infty$-груп, як у~\eqref{eq:normal-fiber-sequence}, то ми отримуємо відповідну спектральну послідовність з $E_2$-сторінкою
\[
  H_{\mathrm{Grp}}^p(G \dblslash K, H_{\mathrm{Grp}}^q(K, A))
\]
і яка збігається до $H_{\mathrm{Grp}}^n(G, A)$, де $A$ — будь-який усічений зв'язний спектр, який міг би бути навіть лівим $G$-модулем, і в цьому випадку ми відтворюємо \emph{спектральну послідовність Хохшильда-Серра}.

\section{Формалізація}
\label{sec:formalization}
Ми формалізували багато результатів цієї статті. Ми використовуємо доказовий асистент Lean 2. Це стара версія доказового асистента Lean (версія 3.3 станом на січень 2018 року). Ми використовуємо стару версію, оскільки нова версія офіційно не підтримує HoTT, хоча існує експериментальна бібліотека для HoTT, але вона не має такої кількості теорій, як бібліотека в Lean 2.

Бібліотека Lean 2 HoTT розділена на дві частини: основну бібліотеку та формалізацію спектральних послідовностей. Ми працювали в останній, щоб мати можливість використовувати результати з цього репозиторію, такі як теореми про простори Ейленберга-Маклейна та точкові відображення. Усі результати цієї статті викладені в одному файлі, хоча для багатьох результатів основна частина доведення знаходиться в іншому місці (в Emacs натисніть на назву та натисніть \texttt{M-.}, щоб знайти визначення).

Щоб зібрати файл, встановіть Lean 2 згідно з інструкціями з цього репозиторію, а потім завантажте репозиторій Spectral і компилюйте його (ви можете використовувати команду \texttt{path/to/lean2/bin/linja} у командному рядку, щоб скомпілювати бібліотеку, в якій ви перебуваєте). Репозиторій Spectral містить деякі недоведені результати, позначені \textbf{sorry}. Ви можете написати \texttt{print axioms theoremname} у файлі, щоб переконатися, що \textbf{sorry} не використовується в доведенні.

\section{Висновок}
\label{sec:conclusion}

Ми представили теорію та формалізацію вищих груп у HoTT і довели, що для структур на рівні множин ми відновлюємо добре відомі об'єкти: групи та абелеві групи. Можливим наступним кроком було б зробити те саме для об'єктів 1-типу. Відповідні алгебраїчні об'єкти мають довгу історію. Суворі 2-групи передують теорії категорій, оскільки вони виникають у дослідженні Вайтхедом \emph{схрещених модулів}~\cite{Whitehead1949}. Теорія слабких 2-груп була розпочата учнем Гротендіка Хоанг Суан Сінем~\cite{Sinh1975} і далі розвинена в~\cite{BaezLauda2004}. У HoTT має бути можливим довести, що слабкі 2-групи та схрещені модулі еквівалентні 2-групам у нашому розумінні, коли ми використовуємо відповідні правильні поняття еквівалентності.

Симетричні 2-групи за теоремою про стабілізацію є такими ж, як 1-усічені симетричні спектри. Вони описуються простіше, ніж довільні схрещені модулі, як \emph{групоїди Пікара}. Це частина стабільної гомотопічної гіпотези~\cite{Gurski-stable-homotopy-hypothesis,JohnsonOsorno}. Також має бути можливим розвинути теорію групоїдів Пікара в HoTT і таким чином довести відповідну стабільну гомотопічну гіпотезу.

Вищі групи інтенсивно вивчалися в теорії гомотопій, зокрема після $p$-поповнення для простого числа $p$. \emph{$p$-компактна група} — це $\bF_p$-локальна $\infty$-група, носій якої є $\bF_p$-скінченним, див.~\cite{DwyerWilkerson1994}. Вони є хорошими гомотопічними аналогами груп Лі і добре взаємодіють з компактними групами Лі, наприклад:

\begin{theorem}[\cite{DwyerZabrodsky1987}]
  Нехай $P$ — $p$-торальна група, а $G$ — компактна група Лі. Тоді $\trunc{BP \to_\pt BG}_0$ ізоморфно класам спряженості гомоморфізмів з $P$ в $G$.
\end{theorem}

Вищі групи також відіграють особливо помітну роль у розвитку квантової теорії поля в когезивній гомотопічній теорії типів~\cite{ShreiberShulman}. У когезивній теорії типів ми справді можемо вловити топологічну або гладку структуру груп та їх класифікуючих типів, а отже, правильно розвинути теорію Лі, включаючи її узагальнення на вищі групи. Усі наші результати використовують лише основну частину HoTT, і тому вони залишаються в силі також у когезивній HoTT.

Зауважте, що ми критично використали трюк для вивчення вищих груп у HoTT, а саме те, що вони можуть бути представлені точковими зв'язними типами. Альтернативою було б визначити їх як групоподібні алгебри для операди маленьких $k$-кубів $E_k$. Але це вимагає саме такої нескінченної вежі умов когерентності, яку ми ще не знаємо, як визначити в HoTT. (Або чи це взагалі можливо.) Таким чином, хоча ми маємо тип вищих груп, ми не маємо типу вищих моноїдів (загальних $E_k$-алгебр). Таким чином, їх теорія та відповідна теорема про стабілізацію наразі недоступні для HoTT.

\section*{Подяка}
  Автори вдячні за підтримку Управлінню наукових досліджень ВПС США через грант MURI FA9550-15-1-0053. Будь-які думки, висновки та рекомендації, висловлені в цьому матеріалі, належать авторам і не обов'язково відображають погляди AFOSR.

\begin{thebibliography}{99}

\bibitem{chtt}
C. Angiuli, R. Harper, and T. Wilson.
\newblock Computational higher-dimensional type theory.
\newblock In {\em Proceedings of the 44th ACM SIGPLAN Symposium on Principles of Programming Languages (POPL 2017)}, pages 680--693. ACM, 2017.

\bibitem{AwodeyWarren2009}
S. Awodey and M. A. Warren.
\newblock Homotopy theoretic models of identity types.
\newblock {\em Math. Proc. Cambridge Philos. Soc.}, 146(1):45--55, 2009.

\bibitem{BaezDolan1998}
J. C. Baez and J. Dolan.
\newblock Categorification.
\newblock In {\em Higher category theory (Evanston, IL, 1997)}, volume 230 of {\em Contemp. Math.}, pages 1--36. Amer. Math. Soc., 1998.

\bibitem{BaezLauda2004}
J. C. Baez and A. D. Lauda.
\newblock Higher-dimensional algebra. V: 2-groups.
\newblock {\em Theory Appl. Categ.}, 12:423--491, 2004.

\bibitem{realprojective}
U. Buchholtz and E. Rijke.
\newblock The real projective spaces in homotopy type theory.
\newblock In {\em 32nd Annual ACM/IEEE Symposium on Logic in Computer Science (LICS 2017)}, pages 1--8. IEEE, 2017.

\bibitem{Camarena}
O. A. Camarena.
\newblock The homotopy fiber of the map on classifying spaces, 2017.
\newblock URL \url{http://www.matem.unam.mx/omar/notes/hofib-grphom.html}.

\bibitem{CCHM2016}
C. Cohen, T. Coquand, S. Huber, and A. Mörtberg.
\newblock Cubical type theory: a constructive interpretation of the univalence axiom.
\newblock In {\em 21st International Conference on Types for Proofs and Programs (TYPES 2015)}, 2016.

\bibitem{SerreSpectralSequence}
F. van Doorn, J. Avigad, S. Awodey, U. Buchholtz, E. Rijke, and M. Shulman.
\newblock Spectral sequences in homotopy type theory, 2018.
\newblock URL \url{https://github.com/cmu-phil/Spectral}.

\bibitem{vDvRB2017HoTTLean}
F. van Doorn, J. von Raumer, and U. Buchholtz.
\newblock Homotopy type theory in {Lean}.
\newblock In {\em Interactive Theorem Proving (ITP 2017)}, pages 479--495. Springer, 2017.

\bibitem{DwyerWilkerson1994}
W. G. Dwyer and C. W. Wilkerson.
\newblock Homotopy fixed-point methods for {Lie} groups and finite loop spaces.
\newblock {\em Ann. Math. (2)}, 139(2):395--442, 1994.

\bibitem{DwyerZabrodsky1987}
W. G. Dwyer and A. Zabrodsky.
\newblock Maps between classifying spaces.
\newblock In {\em Algebraic topology, Barcelona, 1986}, pages 106--119. Springer, 1987.

\bibitem{Elmendorf1983}
A. D. Elmendorf.
\newblock Systems of fixed point sets.
\newblock {\em Trans. Amer. Math. Soc.}, 277:275--284, 1983.

\bibitem{Grothendieck1983}
A. Grothendieck.
\newblock Pursuing stacks.
\newblock Manuscript, 1983.

\bibitem{Gurski-stable-homotopy-hypothesis}
N. Gurski.
\newblock The stable homotopy hypothesis and categorified abelian groups.
\newblock Blog post, The $n$-Category Café, 2018.

\bibitem{FavoniaHarper2016}
K. B. Hou (Favonia) and R. Harper.
\newblock Covering spaces in homotopy type theory.
\newblock In {\em 22nd International Conference on Types for Proofs and Programs (TYPES 2016)}, 2016.

\bibitem{JohnsonOsorno}
N. Johnson and A. M. Osorno.
\newblock Modeling stable one-types.
\newblock {\em Theory Appl. Categ.}, 26(20):520--537, 2012.

\bibitem{FinsterLicata2014}
D. R. Licata and E. Finster.
\newblock {Eilenberg-MacLane} spaces in homotopy type theory.
\newblock In {\em CSL-LICS '14}, pages 66:1--66:9. ACM, 2014.

\bibitem{LurieHTT}
J. Lurie.
\newblock {\em Higher Topos Theory}.
\newblock Princeton University Press, 2009.

\bibitem{Moura2015}
L. de Moura, S. Kong, J. Avigad, F. van Doorn, and J. von Raumer.
\newblock The {Lean} theorem prover (system description).
\newblock In {\em CADE-25}, pages 378--388. Springer, 2015.

\bibitem{joinconstruction}
E. Rijke.
\newblock The join construction, 2017.
\newblock arXiv:1701.07538.

\bibitem{ShreiberShulman}
U. Schreiber and M. Shulman.
\newblock Quantum gauge field theory in cohesive homotopy type theory.
\newblock In {\em Proceedings 9th Workshop on Quantum Physics and Logic}, volume 158 of {\em EPTCS}, pages 109--126, 2014.

\bibitem{Shulman-classifying}
M. Shulman.
\newblock The univalent perspective on classifying spaces.
\newblock Blog post, The $n$-Category Café, 2015.

\bibitem{ShulmanEI}
M. Shulman.
\newblock Univalence for inverse {EI} diagrams.
\newblock {\em Homology, Homotopy and Applications}, 19(2):219--249, 2017.

\bibitem{Sinh1975}
H. X. Sinh.
\newblock {\em Gr-catégories}.
\newblock PhD thesis, Université Paris Diderot (Paris 7), 1975.

\bibitem{TheBook}
The {Univalent Foundations Program}.
\newblock {\em Homotopy Type Theory: Univalent Foundations of Mathematics}.
\newblock Institute for Advanced Study, 2013.

\bibitem{Voevodsky06}
V. Voevodsky.
\newblock A very short note on homotopy $\lambda$-calculus, 2006.

\bibitem{Whitehead1949}
J. H. C. Whitehead.
\newblock Combinatorial homotopy. {II}.
\newblock {\em Bull. Amer. Math. Soc.}, 55:453--496, 1949.

\end{thebibliography}

\end{document}
